
des $G_{\mathbb{R}}$-modules de dim.\ finie sur $\mathbb{R}$. La
\uncertain{condition} qui exprime cette effectivité est \struck{singulière}
par les \struck{réels} noyaux définis par les $v(s')$ [sous forme
« \uncertain{analytique} »]. C'est un fait qui exprime [le fait que les]
Frobenius
\[
  f(s') \in G_{\mathbb{R}}^{\xi(s')}(\mathbb{R})
\]
ont \ill{} par leurs valeurs propres, dans \ill{} représentation \ill{} de
$G_{\mathbb{R}}^{\xi(s')}$ \struck{$(\mathbb{R})$} \ill{} \ill{} de
$G_{\mathbb{C}}^{\xi(s')}$ --- sont des \emph{entiers} algébriques, ce qui
est une condition arithmétique assez draconienne !
\note{« entiers » est souligné. Le premier mot de la deuxième ligne,
« singulière », est souligné puis biffé ; l'exposant de $f$ se lit
$(s')$. Le texte s'arrête au tiers supérieur du feuillet, le reste est
blanc}

\section{Théorie de Galois des Motifs. Généralités}

\note{deux lignes de sa main, au crayon, soulignées, seules sur le feuillet
188, qui sert de couverture ; elles deviennent ici le titre de section}

\subsection*{Foncteurs fibres motiviques et gp fondamental motivique}

\note{titre de sa main, seul sur le feuillet 189, qui sert de couverture}

